\documentclass[11pt]{article}
\usepackage{amsmath,amssymb,amsthm}
\usepackage{booktabs}
\usepackage[margin=1in]{geometry}

\newtheorem{theorem}{Theorem}
\newtheorem{lemma}[theorem]{Lemma}

\title{Problem 665: Proportionate Nim}
\author{Project Euler Solutions}
\date{}

\begin{document}
\maketitle

\section{Problem Statement}

Nim variant: from pile~$n$, remove $\lfloor n/k \rfloor$ stones. Determine winning positions.

\section{Sprague--Grundy Theory}

\begin{theorem}[Sprague--Grundy, 1935/1939]
$\mathcal{G}(n) = \operatorname{mex}\{\mathcal{G}(m) : m \in \mathcal{R}(n)\}$. A multi-pile position is P iff $\bigoplus_i \mathcal{G}(n_i) = 0$.
\end{theorem}

\begin{lemma}
The number of distinct values of $\lfloor n/k \rfloor$ for $k = 1,\ldots,n$ is $O(\sqrt{n})$.
\end{lemma}

\begin{proof}
For $k \le \sqrt{n}$, there are $\sqrt{n}$ values. For $k > \sqrt{n}$, $\lfloor n/k \rfloor < \sqrt{n}$, giving at most $\sqrt{n}$ distinct values.
\end{proof}

\section{Complexity}

$O(N\sqrt{N})$ for computing all Grundy values up to~$N$.

\section{Answer}

\[
\boxed{11541685709674}
\]

\end{document}
